\title{Controlling Text-to-Image Diffusion by Orthogonal Finetuning}

\begin{abstract}
Large-scale text-to-image diffusion models have demonstrated remarkable proficiency in synthesizing photorealistic images from textual descriptions. A significant challenge that remains is devising efficient mechanisms to steer or modulate these advanced models for diverse downstream applications. In response, we propose a systematic finetuning technique, Orthogonal Finetuning~(OFT), tailored to customize text-to-image diffusion models for specific tasks. Distinct from existing strategies, OFT can theoretically maintain hyperspherical energy, which encapsulates the relational structure among neurons when projected onto a unit hypersphere. We identify this preservation as key to retaining the semantic generative potential of text-to-image diffusion models during adaptation. To further enhance the robustness of the finetuning process, we introduce Constrained Orthogonal Finetuning (COFT), which enforces a supplementary radius restriction on the hypersphere. Our investigation concentrates on two critical adaptation scenarios: subject-driven generation—where, with a handful of subject images and a textual prompt, the objective is to produce contextually diverse renditions of the subject—and controllable generation, wherein the model integrates auxiliary control signals. Empirical results indicate that the OFT approach surpasses previous techniques in both image generation fidelity and finetuning convergence rate.
\end{abstract}

\section{Introduction}

Recent advancements in text-to-image diffusion models~\cite{saharia2022photorealistic,ramesh2022hierarchical,rombach2022high} have set new standards for high-quality, text-driven image synthesis. Nevertheless, reliance solely on textual prompts often yields ambiguity and falls short in offering precise, fine-tuned control over image attributes. To overcome these constraints, our work focuses on two pivotal categories of text-to-image generation tasks:

\begin{itemize}[leftmargin=*,nosep]

    \item \textbf{Subject-driven generation}~\cite{ruiz2023dreambooth}: Provided with only several visual examples of a target subject, the aim is to synthesize novel images depicting the same subject in varied contexts, guided by a textual prompt.
    \item \textbf{Controllable generation}~\cite{zhang2023adding,mou2023t2i}: When supplied with extra control inputs (such as canny edge maps or segmentation cues), the model’s goal is to generate images aligned with both the control information and a text prompt.
\end{itemize}

Both tasks fundamentally address the challenge of how to adapt text-to-image diffusion models to new tasks while maintaining the high-quality generative abilities achieved during pretraining. The main criteria for an ideal finetuning approach are as follows: (1) \emph{training efficiency}, defined by the minimization of both trainable parameters and training epochs, and (2) \emph{generalizability preservation}, which refers to sustaining the diversity and fidelity of the generated outputs. Common practices in finetuning involve either modifying existing neuron weights using a small learning rate (\emph{e.g.}, \cite{ruiz2023dreambooth}), or integrating lightweight modules that re-parameterize certain weights (\emph{e.g.}, \cite{hulora2022,zhang2023adding}). However, despite their straightforwardness, these strategies fall short of guaranteeing that the original generative strengths from pretraining are maintained. Moreover, there is a notable gap in theoretical understanding regarding the construction of effective finetuning mechanisms, as well as the optimal selection of hyperparameters, including the duration of training. This difficulty is compounded by the absence of robust metrics for evaluating to what extent the pretrained generative capabilities are retained. Existing techniques generally operate under the assumption that minimizing the Euclidean distance between the weights of the finetuned and pretrained models leads to superior retention of pretrained performance, and consequently they rely on extremely small learning rates. While this rationale may sometimes yield positive results, we contend that focusing solely on Euclidean proximity is insufficient to encapsulate the semantic consistency between models. Thus, employing a more structurally informed metric for assessing the deviation between finetuned and pretrained models could substantially advance both the fidelity of pretraining preservation and the robustness of the finetuning process.

Building on empirical findings that hyperspherical similarity effectively captures semantic relationships~\cite{liu2017deep,liu2018decoupled,chen2020angular}, we leverage hyperspherical energy~\cite{liu2018learning} as a tool to represent the pairwise interactions among neurons. Specifically, hyperspherical energy quantifies the aggregate hyperspherical similarity (such as cosine similarity) between all neuron pairs within a given layer, thereby reflecting how uniformly neurons are distributed on the unit hypersphere~\cite{liu2021learning}. Our working hypothesis is that optimal finetuning should result in minimal deviation of hyperspherical energy from that of the original pretrained model. One straightforward strategy would be to introduce a regularization term to maintain constant hyperspherical energy during finetuning; however, this does not ensure an effectively minimized energy difference. To address this, we exploit a fundamental invariance of hyperspherical energy—namely, the property that applying the same orthogonal transformation across all neurons preserves their pairwise hyperspherical similarity. Motivated by this characteristic, we introduce Orthogonal Finetuning~(OFT), a technique designed to adapt large text-to-image diffusion models to new tasks without altering their hyperspherical energy. The principal concept is to learn a single orthogonal transformation applicable to all neurons within a layer, thus preserving the angles between neuron vectors. Alternatively, OFT can be interpreted as redefining the basis in which neurons are represented within each layer. Moreover, recognizing that a reduced Euclidean distance between the finetuned and pretrained models correlates with better retention of pretraining efficacy, we present an extension: Constrained Orthogonal Finetuning~(COFT), which further restricts the updated model to reside within a fixed-radius hypersphere centered at the pretrained neuron positions.

The rationale underlying the effectiveness of orthogonal transformation for neuron finetuning draws in part from the concept of 2D Fourier transform, where an image is represented by its magnitude and phase spectra. In this decomposition, the phase spectrum—which encapsulates angular relationships between the input and the basis—captures the primary semantic content. Notably, combining the phase spectrum of an image with an arbitrary magnitude spectrum can still result in a reconstructed image that retains the original semantic meaning. This indicates that manipulating the directions of neurons plays a crucial role in semantically altering the generated image, which is essential for both subject-driven and controllable image generation. Nevertheless, unrestricted modifications to neuron directions can significantly disrupt the generative capabilities established during pretraining. To address this, we advocate for maintaining the angles between every pair of neurons, predicated on the assumption that these angular relationships are integral to neural network knowledge representation. From this perspective, it is natural to employ a layer-wise, shared orthogonal transformation for neurons, ensuring the preservation of hyperspherical energy across layers.

Additionally, our methodology is informed by the framework of orthogonal over-parameterized training~\cite{liu2021orthogonal}, wherein classification neural networks are trained from scratch by applying orthogonal transformations to their random initializations. The reason for this is that randomly initialized networks typically exhibit low expected hyperspherical energy, and the objective in~\cite{liu2021orthogonal} is to keep this energy minimal throughout training, since low hyperspherical energy correlates with improved classification generalization~\cite{liu2018learning,lin2020regularizing}. The results in~\cite{liu2021orthogonal} demonstrate that orthogonal transformations possess sufficient expressiveness for training classification networks that generalize well. Our focus, in contrast, is on the finetuning of text-to-image diffusion models to enhance both their controllability and generative power on downstream tasks. The distinction between OFT and~\cite{liu2021orthogonal} is apparent in two key ways. Firstly, while the latter aims to reduce hyperspherical energy, OFT seeks to retain the hyperspherical energy at the level established by the pretrained model, thereby safeguarding the essential structure learned during pretraining—minimizing this energy during diffusion model finetuning could erase valuable semantic structures. Secondly, OFT explicitly minimizes divergence from the pretrained configuration, resulting in a constrained version, unlike~\cite{liu2021orthogonal}, which imposes no such restriction. The principal challenge in finetuning is to strike an optimal balance between adaptability and preservation of learned knowledge, and we contend that our OFT paradigm is particularly well-suited to achieving this. Our main contributions are as follows:

\begin{itemize}[leftmargin=*,nosep]

    \item We introduce a new finetuning technique—Orthogonal Finetuning—for enhancing the controllability of text-to-image diffusion models. To boost robustness, we further propose a constrained form of our method that restricts angular shifts relative to the original pretrained model.
    \item In contrast to previous finetuning strategies, OFT enables model adaptation while formally maintaining the hyperspherical energy. Through empirical evaluation, we identify hyperspherical energy as a critical indicator of a pretrained model’s ability to preserve generative semantics.
    \item The OFT approach is utilized across two scenarios: subject-driven generation and controllable generation. We present extensive empirical evidence showing that our approach leads to notable gains over existing methods in terms of output quality, training speed, and stability during finetuning. Additionally, OFT exhibits superior sample efficiency, successfully converging with markedly fewer training samples and epochs.
    \item For tasks demanding controllable generation, we present a novel control consistency metric to quantify controllability. This involves estimating the control signal from a generated image and contrasting it with the original control input. The results further substantiate OFT’s superior capability in achieving consistent control.
\end{itemize}

\section{Related Work}

\textbf{Text-to-image diffusion models}. Significant advances~\cite{saharia2022photorealistic,ramesh2022hierarchical,rombach2022high,gu2022vector,nichol2022glide} have been achieved in the field of text-to-image synthesis, propelled by breakthroughs in diffusion-based generative frameworks~\cite{ho2020denoising,song2021score,song2021denoising,dhariwal2021diffusion} as well as progress in joint vision-language modeling~\cite{su2020vlbert,lu2019vilbert,radford2021learning,wang2021simvlm,li2022blip,li2023blip,alayrac2022flamingo,schuhmann2022laion}. Notably, GLIDE~\cite{nichol2022glide} and Imagen~\cite{saharia2022photorealistic} focus on pixel-space diffusion; GLIDE optimizes its text encoder alongside the diffusion process using paired data, whereas Imagen leverages a static, pretrained text encoder. Alternatively, Stable Diffusion~\cite{rombach2022high} and DALL-E2~\cite{ramesh2022hierarchical} adopt a latent space modeling approach: Stable Diffusion makes use of VQ-GAN~\cite{esser2021taming} to establish its latent codebook, while DALL-E2 utilizes CLIP~\cite{radford2021learning} for generating a unified embedding space bridging images and text. Aside from diffusion-based models, research into generative adversarial networks~\cite{reed2016generative,zhang2017stackgan,xu2018attngan,li2019controllable} and autoregressive models~\cite{ramesh2021zero,ding2021cogview,wu2022nuwa,yuscaling} has also contributed to advancements in text-conditioned image generation. OFT itself is agnostic to model architecture and can be seamlessly integrated with any text-to-image diffusion system.

\textbf{Subject-driven generation}. Techniques for maintaining subject integrity include approaches such as~\cite{avrahami2022blended,nichol2022glide}, where user-provided masks are incorporated as supplementary constraints. Inversion-based techniques~\cite{choi2021ilvr,dhariwal2021diffusion,ramesh2022hierarchical,gal2022image} are leveraged for altering context while leaving the main subject unaltered. Similarly, \cite{hertz2022prompt} supports both localized and global edits, yet without needing explicit masks. Despite these methods, preserving fine-grained identity characteristics of the subject remains challenging. The Pivotal Tuning framework~\cite{roich2022pivotal} circumvents this by adapting a generator near an initial inversion in the latent space, incorporating regularization to maintain subject identity. In a related manner, \cite{nitzan2022mystyle} constructs a personalized generative prior for facial images using individual datasets. Although \cite{casanova2021instance} permits generation of distinct variations for a given instance, some identity-specific details can be lost. DreamBooth~\cite{ruiz2023dreambooth} achieves subject-specific adaptation through finetuning with a tailored token and a handful of subject photos, employing a combination of reconstruction and class-specific regularization losses. While OFT builds upon the DreamBooth workflow, it departs from simple low-rate finetuning and instead utilizes orthogonal parameter

\textbf{Controllable generation}. Image-to-image translation is fundamentally a controllable generation task. Traditional approaches in this domain predominantly utilize conditional generative adversarial networks~\cite{isola2017image,zhu2017toward,wang2018high,park2019semantic,choi2018stargan}. Diffusion models have also been adopted for this translation process~\cite{saharia2022palette,voynov2022sketch,wang2022pretraining}. In recent developments, ControlNet~\cite{zhang2023adding} has introduced a mechanism for governing a pretrained diffusion model by finetuning it with supplementary control cues, achieving remarkable performance in controllable generation. Similarly, T2I-Adapter~\cite{mou2023t2i} contemporaneously proposes finetuning pretrained diffusion models to enhance control over image generation. Consistent with the experimental frameworks of \cite{zhang2023adding,mou2023t2i}, we implement OFT to adapt pretrained diffusion models, consistently achieving superior controllability even with fewer training samples and reduced finetuning parameters. Importantly, OFT incurs no additional computational burden at inference time.

\textbf{Model finetuning}. The practice of adapting large pretrained models for downstream tasks has gained significant traction in recent years~\cite{devlin2018bert,brown2020language,he2022masked}. Among various finetuning strategies, adaptation-based techniques (\emph{e.g.}, \cite{hulora2022,houlsby2019parameter,pfeiffer2020adapterfusion}) have been extensively explored within the natural language processing community. OFT bears closest resemblance to LoRA~\cite{hulora2022}, which relies on a low-rank structure for weight updates during the finetuning process. Differentiating itself, OFT implements a multiplicative neuron weight update using layer-wise shared orthogonal transformations. This approach mathematically maintains the angular relationships among neurons in each layer, offering increased stability.

\section{Orthogonal Finetuning}

\subsection{Why Does Orthogonal Transformation Make Sense?}

We begin by examining the motivation for using orthogonal transformations in finetuning text-to-image diffusion models. This discussion is divided into two key considerations: (1) the rationale for tuning neuron angles (\emph{i.e.}, directions), and (2) the justification for utilizing orthogonal transformations as the means for adjusting these angles.

In addressing the first question, we build upon empirical findings from \cite{liu2018decoupled,chen2020angular}, which demonstrate that differences in feature angles serve as a meaningful measure of the semantic distance between features. Notably, SphereNet~\cite{liu2017deep} demonstrates that constraining all neuron outputs to reside on a unit hypersphere during training leads to neural networks with similar capacity and, in some cases, superior generalization, suggesting that the directional component of neuron activations effectively encapsulates essential data characteristics. To more concretely illustrate the significance of neuron direction, we present a controlled experiment (Figure~\ref{fig:toy_ae}) where a canonical convolutional autoencoder is trained using flower image data. During training, standard inner products are utilized to generate feature maps ($z$ denotes the output of applying the convolution kernel $\bm{w}$ to input $\bm{x}$ within a sliding window). When evaluating, we test three alternatives for obtaining the feature map: (a) the same inner product as during training, (b) isolating magnitude information, and (c) using only angular (directional) information. The visualizations in Figure~\ref{fig:toy_ae} reveal that preserving only angular features is nearly sufficient for reconstructing the input images, whereas the magnitude alone provides no meaningful content. Importantly, cosine-based activation (c) is never employed during training—the model is always trained using inner product calculations. These observations point to neuron directions as the primary conveyors of the input images' semantic content. Thus, manipulating neuron angles may be a more promising avenue for altering the underlying semantics of the representations.

For the second question, a straightforward method to adjust neuron directions involves globally rotating or reflecting all neurons within a layer, which corresponds to applying orthogonal transformations. While more complex angular modifications—such as rotating each neuron independently—could provide additional adaptability, orthogonal transformations strike a practical balance between flexibility and maintaining structural coherence. Further, as demonstrated in \cite{liu2021orthogonal}, orthogonal transformations suffice to enable expressive learning in neural architectures. To empirically substantiate this claim, we conduct a regularization analysis using the ControlNet~\cite{zhang2023adding} setup, with results presented in Figure~\ref{fig:ortho}. The findings show that our vanilla OFT consistently yields stable training progress and attains precise control by epoch 20. Conversely, relaxing the orthogonality constraint leads to failed image synthesis and loss of controllability. These results highlight the essential regulatory function that orthogonality enforces within OFT.

\subsection{General Framework}

The standard approach to finetuning often involves applying gradient descent with a relatively low learning rate to adapt the weights of a model, or selected layers thereof. By employing such a small learning rate, this method inherently restricts the deviation of the new model from its pretrained counterpart, effectively encouraging minimal alterations and resulting in an optimization objective akin to minimizing $\thickmuskip=2mu \medmuskip=2mu \| \bm{M} - \bm{M}^0\|$, where $\bm{M}$ and $\bm{M}^0$ represent the finetuned and original weights, respectively. While this soft constraint is intuitively reasonable, it may not be sufficiently restrictive for fine-tuning large-scale models. To overcome this limitation, LoRA augments the process by enforcing that the weight update obeys a low-rank restriction, i.e., $\thickmuskip=2mu \medmuskip=2mu \text{rank}(\bm{M} - \bm{M}^0)=r'$ with $r'$ chosen to be small. Contrasting with LoRA’s low-rank regularization, OFT introduces a different type of constraint, focusing on preserving pairwise neuron similarity: $\thickmuskip=2mu \medmuskip=2mu \| \text{HE}(\bm{M})-\text{HE}(\bm{M}^0) \|=0$, where the function $\text{HE}(\cdot)$ captures the hyperspherical energy of the corresponding weight matrix. 

To provide concrete intuition, consider a fully connected layer parameterized by $\thickmuskip=2mu \medmuskip=2mu \bm{W}=\{\bm{w}_1,\cdots,\bm{w}_n\}\in\mathbb{R}^{d\times n}$, in which $\bm{w}_i\in\mathbb{R}^{d}$ stands for the weights of the $i$-th neuron (with pretrained weights $\bm{W}^0$). The output vector, $\thickmuskip=2mu \medmuskip=2mu \bm{z}\in\mathbb{R}^n$, is calculated as $\thickmuskip=2mu \medmuskip=2mu \bm{z}=\bm{W}^\top\bm{x}$, where $\thickmuskip=2mu \medmuskip=2mu \bm{x}\in\mathbb{R}^d$ denotes the input. Under the OFT framework, the goal can be viewed as minimizing the change in hyperspherical energy incurred through finetuning, relative to the original pretrained model:

\begin{equation}\label{eq:oft_principle}
\footnotesize
\min_{\bm{W}} \left\| \text{HE}(\bm{W})-\text{HE}(\bm{W}^0)\right\|~~\Leftrightarrow~~\min_{\bm{W}} \bigg{\|} \sum_{i\neq j} \|\hat{\bm{w}}_i-\hat{\bm{w}}_j\|^{-1}-\sum_{i\neq j} \|\hat{\bm{w}}^0_i-\hat{\bm{w}}^0_j\|^{-1}\bigg{\|}
\end{equation}

In this context, each normalized neuron is given by $\thickmuskip=2mu \medmuskip=2mu \hat{\bm{w}}_i:=\bm{w}_i/\|\bm{w}_i\|$, and the hyperspherical energy for a layer $\bm{W}$ is defined as $\thickmuskip=2mu \medmuskip=2mu \text{HE}(\bm{W}):= \sum_{i\neq j} \|\hat{\bm{w}}_i-\hat{\bm{w}}_j\|^{-1}$. Evidently, the lowest attainable value for Eq.~\eqref{eq:oft_principle} is zero, which occurs whenever $\bm{W}$ results from applying an orthogonal transformation to $\bm{W}^0$ (i.e., $\thickmuskip=2mu \medmuskip=2mu \bm{W}=\bm{R}\bm{W}^0$ with $\thickmuskip=2mu \medmuskip=2mu \bm{R}\in\mathbb{R}^{d\times d}$ being an orthogonal matrix—either a rotation if $\det \bm{R}=

\begin{equation}\label{eq:oft_general}
    \footnotesize
    \bm{z}=\bm{W}^\top\bm{x}=(\bm{R}\cdot\bm{W}^0)^\top\bm{x},~~~\text{s.t.}~\bm{R}^\top\bm{R}=\bm{R}\bm{R}^\top=\bm{I}
\end{equation}

Here, $\bm{W}$ stands for the weight matrix obtained through OFT-finetuning, and $\bm{I}$ refers to the identity matrix. Figure~\ref{fig:block} offers a visual representation of OFT. Mirroring the zero initialization approach used in LoRA, it is essential for OFT to adjust the pretrained model strictly from $\bm{W}^0$. To ensure this starting point, the orthogonal matrix $\bm{R}$ is initially set as the identity matrix, aligning the finetuned and pretrained weights at initialization. To maintain the orthogonality constraint on $\bm{R}$, we leverage differentiable orthogonalization techniques, as discussed in \cite{liu2021orthogonal,lezcano2019cheap}. In subsequent sections, we examine methods to enforce orthogonality efficiently from a computational perspective.

\subsection{Efficient Orthogonal Parameterization}\label{sect:eff_ortho}

While traditional approaches like the Gram-Schmidt process allow differentiable orthogonalization, they are typically too computationally demanding for practical large-scale use~\cite{liu2021orthogonal}. To achieve greater efficiency, we make use of the Cayley parameterization to produce the orthogonal matrix. This method involves expressing the orthogonal matrix as $\thickmuskip=2mu \medmuskip=2mu \bm{R}=(\bm{I}+\bm{Q})(I-\bm{Q})^{-1}$, where the matrix $\bm{Q}$ is skew-symmetric, fulfilling $\thickmuskip=2mu \medmuskip=2mu \bm{Q}=-\bm{Q}^\top$. The use of this technique introduces a minor constraint: Cayley parameterization generates only orthogonal matrices with determinant $1$, thus restricting $\bm{R}$ to the special orthogonal group. Empirically, we observe that this restriction has negligible impact on performance. However, parameterizing $\bm{R}$ via the Cayley transform can still incur significant parameter cost for large $d$. To reduce this burden, we suggest structuring $\bm{R}$ as a block-diagonal matrix composed of $r$ separate blocks, yielding the following configuration:

\begin{equation}
    \footnotesize
    \bm{R}=\text{diag}(\bm{R}_1,\bm{R}_2,\cdots,\bm{R}_r)=\begin{bmatrix}
    \bm{R}_1\in \text{O}(\frac{d}{r}) &  &\\
    &\ddots & \\
    & & \bm{R}_r\in \text{O}(\frac{d}{r})
    \end{bmatrix}\in \text{O}(d)
\end{equation}

Here, $\text{O}(d)$ represents the group of $d$-dimensional orthogonal matrices, where $\thickmuskip=2mu \medmuskip=2mu \bm{R} \in \mathbb{R}^{d\times d}$ and, for all $i$, $\thickmuskip=2mu \medmuskip=2mu \bm{R}_i \in \mathbb{R}^{d/r\times d/r}$ denote orthogonal matrices. Setting $\thickmuskip=2mu \medmuskip=2mu r=1$ recovers a conventional, fully unconstrained orthogonal matrix. A $d \times d$ orthogonal matrix is parameterized by $\thickmuskip=2mu \medmuskip=2mu d(d-1)/2$ variables, yielding a parameter complexity of $\mathcal{O}(d^2)$. In contrast, for an $r$-block diagonal orthogonal matrix, there are $\thickmuskip=2mu \medmuskip=2mu d(d/r-1)/2$ parameters, thus achieving a parameter complexity of $\mathcal{O}(d^2/r)$. Additionally, by enforcing parameter sharing among the block matrices, i.e., $\thickmuskip=2mu \medmuskip=2mu\bm{R}_i = \bm{R}_j, \forall i\neq j$, the number of independent parameters can be further reduced, yielding a parameter complexity of $\mathcal{O}(d^2/r^2)$. Notably, all such modifications retain the orthogonality of the matrix $\bm{R}$, and thus preserve the property of maintaining hyperspherical energy.

To compare parameter efficiency between OFT and LoRA, note that LoRA with a low-rank setting $r'$ employs $\thickmuskip=2mu \medmuskip=2mu r'(d+n)$ tunable parameters. When both $r$ and $r'$ scale with the neuron dimensionality $d$ (e.g., $\thickmuskip=2mu \medmuskip=2mu r=r'=\alpha d$, where $0<\alpha\leq 1$), the parameter count for LoRA increases to $\mathcal{O}(d^2+dn)$, whereas for OFT it reduces to $\mathcal{O}(d)$. As a concrete example to highlight this disparity, consider a weight matrix of size $\thickmuskip=2mu \medmuskip=2mu 128\times128$: with $\thickmuskip=2mu \medmuskip=2mu r'=8$, LoRA requires $2,048$ tunable parameters, whereas with $\thickmuskip=2mu \medmuskip=2mu r=8$ and no block sharing, OFT involves $960$ trainable parameters.

\subsection{Constrained Orthogonal Finetuning}

Original OFT can be made more restrictive by confining the adapted model to remain close to the pretrained parameter space. The Constrained Orthogonal Finetuning (COFT) method enforces this using the following computation in the forward pass:

\begin{equation}\label{eq:coft_general}
    \footnotesize
    \bm{z}=\bm{W}^\top\bm{x}=(\bm{R}\cdot\bm{W}^0)^\top\bm{x},~~~\text{s.t.}~\bm{R}^\top\bm{R}=\bm{R}\bm{R}^\top=\bm{I},~\norm{\bm{R}-\bm{I}}\leq \epsilon
\end{equation}

This setup imposes an orthogonality requirement alongside an $\epsilon$-neighborhood constraint with respect to the identity matrix. The Cayley parameterization, detailed in Section~\ref{sect:eff_ortho}, ensures the orthogonality condition. Nonetheless, incorporating the $\epsilon$-constraint into a Cayley-parameterized orthogonal matrix is challenging. To facilitate a better understanding of the Cayley transform, we employ the Neumann series expansion to approximate $\thickmuskip=2mu \medmuskip=2mu \bm{R} = (\bm{I} + \bm{Q})(I-\bm{Q})^{-1}$ as $\thickmuskip=2mu \medmuskip=2mu \bm{R}\approx\bm{I}+2\bm{Q}+\mathcal{O}(\bm{Q}^2)$, given that the Ne

Given that the series converges in operator norm, it follows that the constraint $\thickmuskip=2mu \medmuskip=2mu\|\bm{R}-\bm{I}\|\leq\epsilon$ can be equivalently re-expressed within the Cayley transform. In this formulation, the requirement becomes $\thickmuskip=2mu \medmuskip=2mu \|\bm{Q}-\bm{0}\|\leq\epsilon'$ where $\bm{0}$ denotes a zero matrix and $\epsilon'$ is a distinct tolerance parameter. This newly defined constraint on $\bm{Q}$ is straightforward to impose via projected gradient descent. To initialize the orthogonal matrix $\bm{R}$ to the identity, we set $\bm{Q}$ initially as a zero matrix. The COFT approach may thus be interpreted as the intersection of two explicit conditions: enforcing a minimal change in hyperspherical energy and restricting departure from the pretrained parameter values. While the latter constraint is typically an implicit assumption in most finetuning techniques, COFT makes it explicit. Although OFT already demonstrates strong results, COFT achieves even more stable finetuning dynamics, likely as a consequence of this explicit deviation restriction. As illustrated in Figure~\ref{fig:eps_coft}, the hyperparameter $\epsilon$ plays a central role in determining COFT's behavior. A higher $\epsilon$ allows greater finetuning flexibility, and as $\epsilon$ increases, the model performance approaches that of the OFT variant. Conversely, smaller $\epsilon$ increasingly aligns the finetuned model's behavior with that of the original pretrained text-to-image diffusion model.

\subsection{Re-scaled Orthogonal Finetuning}

We introduce an extension to the OFT framework by learning an additional per-neuron scaling parameter. The rationale behind this is that rescaling individual neuron activations does not affect hyperspherical energy, since any change in magnitude is normalized away. Concretely, our modified forward computation is: $\thickmuskip=2mu \medmuskip=2mu\bm{z}=(\bm{R}\bm{W}^0\bm{D})^\top\bm{x}$\footnote{\textbf{Errata}: The NeurIPS camera-ready version incorrectly reported the forward computation for re-scaled OFT as $\thickmuskip=2mu \medmuskip=2mu\bm{z}=(\bm{D}\bm{R}\bm{W}^0)^\top\bm{x}$. The correct implementation, which was used in all experiments (including those reported in Appendix~\ref{app:rescaled}), is given above.} where $\thickmuskip=2mu \medmuskip=2mu \bm{D}=\text{diag}(s_1,\cdots,s_n)\in\mathbb{R}^{n\times n}$, a diagonal matrix of learnable parameters $s_1,\ldots,s_n$ constrained to be positive. Unlike the standard OFT setup, in which only $\bm{R}$ is optimized (see Eq.~\eqref{eq:oft_general}), our extension introduces both the orthogonal $\bm{R}$ and the diagonal scaling matrix $\bm{D}$ as learnable quantities. This variant—re-scaled OFT—enhances OFT's expressiveness, yet adds only a negligible number of additional parameters. In our main experimental section, we employ the unmodified OFT in order to better isolate and demonstrate the contribution of orthogonal transformation. Nonetheless, we observe in practice that the re-scaled OFT typically outperforms the baseline, as documented in Appendix~\ref{app:rescaled}.

\section{Intriguing Insights and Discussions}

\textbf{OFT applies broadly across network architectures.} OFT, by its design, is adaptable to a variety of neural network types. In the context of Transformer architectures, LoRA is commonly applied to attention layers~\cite{hulora2022}; therefore, for parity, our experimental protocols restrict OFT application to the attention weights. Beyond dense (fully connected) layers, OFT is also naturally extendable to convolutional layers, as its block-diagonal $\bm{R}$ admits interpretation especially well-suited to convolutional networks—a property not shared by LoRA. By setting the number of blocks in $\bm{R}$ to the number of input channels, each block acts as a transformation specific to a single neuron channel, drawing an analogy to depth-wise convolutions~\cite{chollet2017xception}. If all blocks are tied, the entire set of neuron activations undergoes a shared orthogonal transformation across channels. Preliminary empirical evaluations on the application of OFT to convolutional layers are provided in Appendix~

\textbf{Connection to LoRA}. LoRA mitigates abrupt changes to pretrained weights by introducing an additional low-rank component. In contrast, OFT regulates the transformation applied to the pretrained matrix by enforcing orthogonality (thus maintaining full rank), thereby preserving the information learned during pretraining. The forward computation in OFT can be recast as $\thickmuskip=2mu \medmuskip=2mu \bm{z}=(\bm{R}\bm{W}^0)^\top\bm{x}=(\bm{W}^0+(\bm{R}-\bm{I})\bm{W}^0)^\top\bm{x}$, with $\thickmuskip=2mu \medmuskip=2mu (\bm{R}-\bm{I})\bm{W}^0$ serving a role similar to the low-rank adjustment in LoRA. Given that $\bm{W}^0$ is typically of full rank, OFT also effectively executes a low-rank update if $\thickmuskip=2mu \medmuskip=2mu \bm{R}-\bm{I}$ is constrained to be low-rank. Whereas LoRA employs a rank hyperparameter $r'$ to regulate parameter efficiency, OFT uses a diagonal block parameter $r$ to achieve a comparable reduction in trainable parameters. Notably, the two methods exemplify different strategies for parameter efficiency: LoRA leverages low-rankness to decrease parameter counts, while OFT benefits from enforcing a block-diagonal sparsity pattern along with orthogonality.

\textbf{Why OFT converges faster}. As illustrated in Figure~\ref{fig:toy_ae}, the most efficient modification for altering semantic content is achieved by rotating neuron orientations, which is precisely the mechanism OFT employs. Furthermore, the optimization process in OFT can be interpreted as adjusting neuron parameters over a smooth hyperspherical manifold, leading to a more favorable optimization landscape—an observation further corroborated by empirical results in \cite{liu2021orthogonal}.

\textbf{Why not minimize hyperspherical energy}. Our approach notably diverges from \cite{liu2021orthogonal} in that we do not seek to minimize hyperspherical energy. For classification settings, maintaining a set of non-redundant neurons is preferable. Minimizing hyperspherical energy leads to a configuration in which neurons are maximally equidistant on the sphere, which, although beneficial in some contexts, is not a suitable objective in finetuning; such a process risks erasing the valuable structure acquired during pretraining.

\textbf{Balancing flexibility and regularization in finetuning}. Our investigation reveals a fundamental balance between model flexibility and the degree of imposed regularity. Conventional finetuning maximizes flexibility but at the expense of stability, often resulting in model collapse. Despite its simplicity, OFT achieves an effective compromise, maintaining flexibility while introducing regularity through the preservation of pairwise neuron angles. The block-diagonal form serves to further constrain the orthogonal transformation, providing a more stringent regularization mechanism.

\textbf{No inference-time computational cost}. In contrast to approaches such as ControlNet, the OFT framework does not increase the computational burden during inference. Once training is complete, the orthogonal matrix $\bm{R}$ can be merged with the pretrained weight matrix $\bm{W}^0$, resulting in the updated weights $\thickmuskip=2mu \medmuskip=2mu \bm{W}=\bm{R}\bm{W}^0$. This procedure leaves inference throughput unchanged relative to the original pretrained model.

\section{Experiments and Results}

\textbf{Experimental configuration}. For our studies, we employ Stable Diffusion v1.5~\cite{rombach2022high} as the foundational text-to-image model. To maintain equitable conditions, samples for evaluation are selected randomly across all methods. In subject-driven tasks, experimental settings are largely adapted from DreamBooth~\cite{ruiz2023dreambooth}, whereas settings for controllable generation draw primarily from ControlNet~\cite{zhang2023adding} and T2I-Adapter~\cite{mou2023t2i}. To fairly assess against LoRA, OFT or COFT is only applied at the same network layers as LoRA. Additional experimental procedures and comprehensive results are described in Appendix~\ref{app:settings}.

\subsection{Subject-driven Generation}

\textbf{Configuration}. DreamBooth~\cite{ruiz2023dreambooth} and LoRA~\cite{hulora2022} are employed as baseline comparisons. Each approach utilizes the loss objective defined by DreamBooth. For DreamBooth and LoRA, we follow the original methodologies and employ their respective optimal hyperparameters. Expanded experimental outcomes can be found in Appendix~\ref{app:settings}, \ref{app:coft_oft}, \ref{app:more_qualitative}, and \ref{app:failure}.

\textbf{Finetuning stability and convergence}. We begin by analyzing both the stability and speed of convergence during finetuning for DreamBooth, LoRA, OFT, and COFT, with the experimental outcomes illustrated in Figure~\ref{fig:teaser} and Figure~\ref{fig:db_convergence}. The results indicate that COFT and OFT enable stable finetuning of the diffusion model. Within 400 iterations, DreamBooth and OFT variants provide satisfactory control over generation, whereas LoRA does not maintain subject identity effectively. By 2000 iterations, DreamBooth starts producing degenerate outputs and LoRA is no longer able to reliably generate a yellow shirt, instead mistakenly creating yellow fur. In contrast, OFT and COFT continue to deliver robust and predictable image control at this stage. These findings underscore the rapid yet reliable convergence of the OFT method for subject-driven synthesis. Notably, the reduced sensitivity to the finetuning iteration count is beneficial, as it circumvents the necessity for extensive iteration number tuning across various subjects. With OFT and COFT, a relatively high iteration count may be set without meticulous adjustment. Furthermore, for COFT using an appropriate $\epsilon$, both the learning rate and number of iterations can be conveniently chosen without intensive parameter tuning.

\textbf{Quantitative comparison}. Adopting the protocol in \cite{ruiz2023dreambooth}, we perform a quantitative assessment measuring subject fidelity (DINO~\cite{caron2021emerging}, CLIP-I~\cite{radford2021learning}), text prompt alignment (CLIP-T~\cite{radford2021learning}), and the diversity of generated samples (LPIPS~\cite{zhang2018unreasonable}). The CLIP-I metric calculates the mean pairwise cosine similarity between the CLIP embeddings of generated and corresponding real images. DINO functions analogously, though it utilizes ViT S/16 DINO embeddings in place of CLIP representations. CLIP-T represents the mean cosine similarity between text prompt embeddings and those of generated images, as determined by CLIP. Additionally, LPIPS measures the average cosine similarity of generated images sharing the same subject and prompt. According to Table~\ref{table:dreambooth_final}, COFT and OFT significantly exceed DreamBooth and LoRA on the DINO and CLIP-I metrics, and display marginally superior or equivalent results for prompt fidelity and diversity. To mitigate stochastic variation, each method is evaluated by repeatedly finetuning an identical pretrained model with 30 distinct random seeds. Overall, the outcomes confirm that OFT not only provides improved convergence characteristics and stability but also consistently achieves higher final performance compared to the baselines.

\textbf{Qualitative comparison}. To facilitate a clearer appreciation of the advantages offered by OFT, we present several randomly selected examples of subject-driven generation in Figure~\ref{fig:dreambooth_final}. For objectivity, all examples are generated using the same finetuned model per method, and no additional prompt engineering is performed for specific outputs. The variant of each method displayed corresponds to the one achieving the highest validation CLIP metrics. According to Figure~\ref{fig:dreambooth_final}, both OFT and COFT demonstrate strong capabilities in retaining semantic attributes of the subject, whereas LoRA frequently fails to maintain the subject's identity (\emph{e.g.}, in the bowl scenario, LoRA does not preserve subject identity at all). Additionally, OFT and COFT exhibit superior control via text prompts compared to DreamBooth, which, despite effectively preserving the subject, often struggles to ensure the output follows the input prompt (\emph{e.g.}, the initial row of the bowl case). This comparison highlights that the OFT approach provides superior subject fidelity alongside improved controllability. Furthermore, OFT's effectiveness does not heavily depend on the precise tuning of the number of iterations, enabling robust results even as the iteration count increases—an advantage not observed with DreamBooth or LoRA. Further illustrative examples appear in Appendix~\ref{app:more_qualitative}. In addition, a human evaluation reported in Appendix~\ref{app:human} provides further support for the superior performance of OFT.

\subsection{Controllable Generation}

\textbf{Settings}. For benchmarks, we adopt ControlNet~\cite{zhang2023adding}, T2I-Adapter~\cite{mou2023t2i}, and LoRA~\cite{hulora2022} as comparison baselines. We explore three complex controllable generation tasks: transforming Canny edges to images~(C2I) using the COCO dataset~\cite{lin2014microsoft}, segmentation map to image translation~(S2I) on the ADE20K dataset~\cite{zhou2017scene}, and landmark-to-face synthesis~(L2F) with the CelebA-HQ dataset~\cite{karras2018progressive,xia2021tedigan}. Each method finetunes Stable Diffusion~(SD) v1.5 for 20 epochs across all three datasets. Supplementary results can be found in Appendix~\ref{app:more_qualitative},~\ref{app:more_control_task},~\ref{app:failure}.

\textbf{Convergence}. To assess convergence rates, we compare ControlNet, T2I-Adapter, LoRA, and COFT on the L2F task, utilizing both quantitative and qualitative analyses. For quantitative evaluation, we measure the average $\ell_2$ distance between the control face landmarks and the predicted landmarks. Figure~\ref{fig:face_convergence} illustrates how the face landmark error evolves as each model is finetuned over increasing epochs. The results indicate that COFT and OFT both attain substantially faster convergence compared to the other baselines. For instance, LoRA requires 20 epochs to reach the same level of performance that OFT achieves after only 8 epochs. It is important to highlight that OFT and COFT utilize a comparable quantity of trainable parameters to LoRA—far fewer than ControlNet—but converge much more rapidly than previous methods. Additionally, OFT’s accelerated convergence is further corroborated by Figure~\ref{fig:teaser}: the rightmost example demonstrates that OFT achieves strong performance even when trained with only 5\% of the ADE20K dataset, underscoring its superior data efficiency compared to ControlNet and LoRA. In qualitative comparisons, we emphasize OFT, COFT, and ControlNet since ControlNet's landmark errors are closest to ours. As shown in Figure~\ref{fig:face_convergence_qual}, OFT and COFT not only converge consistently but also increasingly align the generated face pose to the target landmarks over training. In contrast, ControlNet exhibits a delayed emergence of controllability, manifesting abruptly after the 8-th epoch, in agreement with the sudden convergence reported in \cite{zhang2023adding}. The stability of convergence in OFT therefore presents a practical advantage: OFT's training process is much smoother and more reliable, which simplifies the search for optimal hyperparameters.

\textbf{Quantitative comparison}. To assess the effectiveness of controllable generation, we propose a control consistency metric. This metric operates by deriving the control signal from the produced image and comparing it to the original control signal provided as input. For the C2I task, Intersection over Union (IoU) and F1 score are used as evaluation metrics. For the S2I task, we report mean IoU as well as mean and overall accuracy. The L2F task is evaluated using the average $\ell_2$ distance between the intended control landmarks and those predicted by the model. Further information about these consistency metrics can be found in Appendix~\ref{app:settings}. Each method is tested under its most favorable hyperparameter settings. The results, summarized in Table~\ref{table:control_gen}, demonstrate that both OFT and COFT deliver substantially stronger and more precise control compared to competing methods. Adapter-based models (such as T2I-Adapter and ControlNet) exhibit slower convergence and inferior final performance. In particular, LoRA shows superior performance to ControlNet on the S2I task, but falls short in the C2I and L2F settings. Overall, we note that LoRA reaches its performance limits even after extensive hyperparameter tuning. In contrast, OFT continues to improve as the trainable parameter count increases, suggesting its potential for further gains. Our quantitative approach to evaluating controllable generation represents a novel contribution of this work, as it enables robust assessment of control fidelity for fine-tuned models—limited primarily by the accuracy of the pre-existing segmentation or detection models employed.

\textbf{Qualitative comparison}. We further conduct a qualitative assessment of OFT and COFT against leading contemporary approaches such as ControlNet, T2I-Adapter, and LoRA. The randomly selected samples displayed in Figure~\ref{fig:control_final} demonstrate that OFT and COFT deliver not only highly realistic and detailed images, but also allow for precise control over the generation process. For the S2I task, it is evident that LoRA fails to produce outputs aligned with the given segmentation maps, whereas ControlNet, OFT, and COFT exhibit strong controllability in image generation. OFT and COFT, relative to ControlNet, are capable of synthesizing images with finer details and more accurate geometry, despite utilizing considerably fewer model parameters. In the C2I scenario, both OFT and COFT can produce semantically consistent images from coarse Canny edge inputs, a challenge where T2I-Adapter and LoRA show significant deficiencies. Regarding the L2F task, our approach excels in preserving correct pose guidance for generated faces, even under complex pose variations. Across all three controlled generation tasks, OFT and COFT consistently outperform current state-of-the-art methods in image quality, validating the advantage of the proposed OFT architecture for controllable image synthesis. For a broader set of qualitative results, Appendix~\ref{app:control_exp} includes additional samples for each task. Furthermore, Appendix~\ref{app:more_control_task} illustrates that OFT extends well to other control applications, including dense pose-to-human body, sketch-to-image, and depth-to-image tasks.

\section{Concluding Remarks and Open Problems}

Building on the insight that angular relationships among neurons are central to shaping visual semantics, we introduce a straightforward yet robust finetuning strategy—orthogonal finetuning—for enhanced control of text-to-image diffusion models. Our approach is applied to two main tasks: subject-driven generation and controllable generation. In comparison with current techniques, OFT offers improved control and stability during finetuning, requiring fewer tunable parameters. Crucially, OFT avoids incurring extra inference costs, making it practical for deployment.

The introduction of OFT also raises several notable open questions. Firstly, OFT enforces orthogonality through Cayley parametrization, necessitating a matrix inversion operation, which poses a constraint on OFT's scalability. While we alleviate this bottleneck by employing block diagonal parametrization, efficiently accelerating this inversion in a differentiable manner is still unresolved. Secondly, OFT possesses a distinctive capability for compositionality: orthogonal matrices derived from separate OFT finetuning sessions can be combined through multiplication and the result remains orthogonal. Investigating whether these orthogonal matrices continue to encapsulate knowledge from all respective downstream tasks is an intriguing avenue for further exploration. Lastly, OFT's parameter efficiency relies substantially on the block diagonal architecture, which unavoidably introduces certain biases and constrains flexibility. Developing approaches to achieve better parameter efficiency with reduced bias and increased flexibility stands as a significant open problem.